\section{Future Work}\label{sec:future-work}

The most immediate next step is a legally and ethically defensible access format. The current public material consists of methodology and collection code only. A fuller corpus release would make the resource usable without re-running the entire pipeline, which matters for researchers without ATProto collection experience, but it should wait until the redistribution questions discussed in Section~\ref{sec:limitations} have been addressed. Possible formats include stable identifiers, strongly pseudonymised text, or controlled access rather than an unrestricted full-text dump.

A further step would be structured community consultation. Because this corpus captures a small and identifiable public, technically valid release formats may still differ in how acceptable they feel to the users whose posts are included. Future work could therefore combine legal and methodological analysis with lightweight community consultation on corpus access, quotation, and redistribution practices.

Adding linguistic annotation would extend the corpus considerably. A candidate pipeline has been prepared: CLASSLA for sentence segmentation, tokenisation, lemmatisation, and Slovene morphosyntactic tags \parencite{ljubesic2019neural}, with Trankit providing Universal Dependencies annotation \parencite{nguyen2021trankit}. This follows the broad direction established by the JANES project for Slovene user-generated content \parencite{erjavec2019janes}. The pipeline has not yet been run over the full corpus due to computational cost.

The corpus can also support follow-up work on non-standard Slovene processing, code-switching between Slovene and neighbouring languages, thread-level dialogue, and community formation during platform migration. The high reply rate (68.3\%) makes thread reconstruction especially relevant, while the temporal structure makes the corpus useful for longitudinal study.

Finally, the collection infrastructure can be extended longitudinally and applied to other small-language communities on Bluesky. That would make it possible to compare how under-resourced languages establish themselves on an open-protocol platform, and to test whether the Slovene pattern reported here is exceptional or part of a broader ATProto trajectory.
